\chapter{Austerity, Scarcity, and Fiscal Morality}
\label{ch:austerity}

\epigraph{The national budget must be balanced. The public debt must be reduced. The arrogance of the authorities must be moderated and controlled.}{Attributed to Cicero (apocryphal)}

\noindent
In 2012, a retired Greek schoolteacher named Dimitris Christoulas walked to Syntagma Square in central Athens and shot himself.
His suicide note read: ``I find no other solution than a dignified end before I start searching through garbage for food.''
His pension had been cut by 40\% under the austerity program imposed by Greece's creditors.
He was 77 years old, had worked his entire adult life, and had committed no financial crime.
The narrative that justified his impoverishment---that Greece had ``lived beyond its means'' and must now ``tighten its belt''---is the subject of this chapter.

Austerity narratives---the claim that public spending must be cut, that governments must ``live within their means,'' that debt is inherently dangerous---are among the most successful doxonomic products of the past century.
They have shaped policy in dozens of countries, redirected trillions of dollars from public services to debt repayment, and helped justify a redistribution of risk from capital to labor.
This chapter analyzes their construction, the institutional infrastructure that produces them, the formal dynamics of fiscal belief, and the distributional consequences of treating a policy choice as a moral imperative.

The doxonomic interest is not primarily whether austerity is good or bad economics (though the evidence is relevant) but how a specific fiscal framework---one whose empirical support is contested and whose distributional consequences are highly unequal---achieved the status of common sense across multiple countries within a few years.

\section{The Household Metaphor}
\label{sec:ch25-household}

\subsection{Definition and mechanism}

\begin{definition}[Fiscal Morality Narrative]
\label{def:fiscal-morality}
\index{fiscal morality}
The \emph{fiscal morality narrative} is the claim that sovereign budgets are analogous to household budgets: governments, like families, must not spend more than they earn, must avoid debt, and must ``tighten their belts'' in hard times.
\end{definition}

This analogy is economically misleading---sovereign currency issuers face fundamentally different constraints than households (they cannot go bankrupt in their own currency, they can create money, and their spending creates income for the private sector)---but doxonomically powerful because it maps complex macroeconomics onto intuitive domestic experience.
Every adult has managed a household budget; almost no one has managed a sovereign treasury.
The analogy exploits this asymmetry in lived experience.

\subsection{The moral dimension}

The household metaphor is not merely an analogy; it is a \emph{moral} claim.
Debt is framed as sin, spending as temptation, austerity as virtue.
\textcite{graeber2011} traces this moral loading of debt across five millennia, showing that the equation of debt with moral failing is a culturally constructed norm that serves the interests of creditors.

\begin{proposition}[Debt-Morality Fusion]
\label{prop:debt-morality}
\index{debt morality}
The fiscal morality narrative fuses two distinct questions:
\begin{enumerate}[nosep]
  \item \textbf{Technical question}: what is the optimal level of public debt given growth, interest rates, and multiplier effects?
  \item \textbf{Moral question}: is it ``right'' for governments to borrow?
\end{enumerate}
By fusing these, the narrative transforms a technical debate (where reasonable economists disagree) into a moral imperative (where ``responsible'' people agree).
This fusion makes the narrative resistant to empirical refutation: even if austerity produces bad economic outcomes, it can still be defended as ``the right thing to do.''
\end{proposition}

\subsection{Institutional amplification}

The household metaphor circulates through:
\begin{itemize}[nosep]
  \item \textbf{Political speech}: ``We have maxed out our credit card'' (repeated by politicians across the spectrum in multiple countries).
  \item \textbf{Media framing}: news reports consistently use household-budget language for sovereign finances---``the government's deficit,'' ``the nation's credit card,'' ``living beyond our means.''
  \item \textbf{Think-tank output}: organizations like the Peterson Foundation (US), the Taxpayers' Alliance (UK), and the Initiative for a Responsible Federal Budget produce materials that frame public debt as a crisis requiring immediate belt-tightening.
  \item \textbf{Popular economics}: bestselling books (e.g., \emph{The Coming Generational Storm}, \emph{The Debt Bomb}) translate the narrative for mass audiences.
\end{itemize}

\section{The Doxonomic Genealogy of Austerity}
\label{sec:ch25-genealogy}

\subsection{The classical foundation}

The austerity narrative draws on a long intellectual tradition:
\begin{itemize}[nosep]
  \item \textbf{Ricardo's equivalence} (1817): the proposition that government borrowing is equivalent to taxation because rational agents anticipate future taxes. In its strong form (which Ricardo himself doubted), it implies that fiscal policy is irrelevant.
  \item \textbf{The ``Treasury view''} (1920s--30s): the British Treasury's doctrine that government spending necessarily ``crowds out'' private investment, implying that stimulus cannot increase total output.
  \item \textbf{Monetarism} (1960s--70s): Friedman's argument that fiscal policy is ineffective and that only monetary policy matters, providing intellectual cover for reducing government's fiscal role.
  \item \textbf{Public choice theory} (1970s--80s): Buchanan and Tullock's argument that government spending reflects bureaucratic self-interest rather than public welfare, implying that less spending is generally better.
\end{itemize}

Each intellectual precursor served a specific institutional function: Ricardo's equivalence was developed in the context of debates about war financing; the Treasury view defended the gold standard; monetarism justified reducing the post-war welfare state; public choice theory provided the theoretical foundation for Thatcher/Reagan spending cuts.

\subsection{The institutional build-out}

The modern austerity infrastructure was built during the neoliberal period (1970s--2000s):

\begin{enumerate}
  \item \textbf{Think tanks}: Heritage Foundation's ``Mandate for Leadership'' (1981) included detailed spending-cut proposals for every federal department. The Cato Institute's annual ``Budget Pork'' reports framed all government spending as waste. In the UK, the Institute of Economic Affairs and the Adam Smith Institute played analogous roles.

  \item \textbf{International institutions}: the IMF's ``structural adjustment programs'' (1980s--90s) conditioned loans to developing countries on austerity measures, creating a global institutional norm that debt reduction trumps social spending.
  The ``Washington Consensus'' codified this norm: fiscal discipline was listed as the first of ten policy prescriptions.

  \item \textbf{Fiscal rules and institutions}: balanced-budget amendments, debt brakes (Germany's \emph{Schuldenbremse}), the Maastricht Treaty's 3\%/60\% rules for eurozone members---all encoded the austerity narrative in law, removing fiscal decisions from democratic deliberation.

  \item \textbf{Credit rating agencies}: Moody's, S\&P, and Fitch evaluate sovereign creditworthiness, creating a private-sector enforcement mechanism for fiscal discipline. A credit downgrade raises borrowing costs, punishing governments that reject the austerity narrative.
\end{enumerate}

\section{Academic Credibility Laundering}
\label{sec:ch25-credibility}

\subsection{The Reinhart-Rogoff episode}

The most vivid example of academic credibility laundering in the austerity narrative is the Reinhart-Rogoff affair.

\begin{example}
\label{ex:reinhart-rogoff}
In January 2010, Carmen Reinhart and Kenneth Rogoff published ``Growth in a Time of Debt,'' claiming that countries with public debt exceeding 90\% of GDP experienced dramatically lower growth.
The paper was not peer-reviewed but was presented at the American Economic Association conference and published as an NBER working paper.

\paragraph{The doxonomic chain.}
\begin{enumerate}[nosep]
  \item \emph{Production}: two Harvard economists with high epistemic capital produced a paper with a clear, quotable result (``90\% is the threshold'').
  \item \emph{Amplification}: the Peterson Foundation promoted the finding; EU Commissioner Olli Rehn cited it in official policy documents; Paul Ryan cited it in his US budget proposal; George Osborne referenced it in justifying UK austerity.
  \item \emph{Capture}: within months, ``the 90\% threshold'' became common sense in policy debates---a number that could be cited to end any argument about whether austerity was necessary.
\end{enumerate}

\paragraph{The collapse.}
In 2013, Thomas Herndon, a graduate student at UMass Amherst, discovered that the paper contained a spreadsheet coding error (five countries were accidentally excluded from a key calculation), selective data exclusion, and an unconventional weighting methodology.
Correcting these errors eliminated the sharp threshold: countries above 90\% debt had lower but still positive growth, and the relationship was far from the cliff-edge the original paper suggested.

\paragraph{The doxonomic aftermath.}
The paper had already served its function: austerity policies justified partly by its findings had been implemented across Europe, affecting millions of people.
The correction received far less institutional amplification than the original claim---a classic asymmetry in doxonomic production, where the initial framing captures common sense and corrections arrive too late to undo the policy consequences.
\end{example}

\subsection{The general pattern}

The Reinhart-Rogoff case illustrates a general doxonomic pattern:

\begin{model}[Academic Credibility Laundering]
\label{mod:credibility-laundering}
\index{credibility laundering!academic}
\begin{enumerate}[nosep]
  \item A researcher at a high-credibility institution produces a finding that supports a pre-existing policy preference.
  \item The finding is amplified through think tanks, media, and policy networks before full peer review.
  \item The finding's epistemic capital (derived from the researcher's institutional affiliation) exceeds its evidential warrant.
  \item Policy decisions are made based on the amplified finding.
  \item When the finding is later corrected or qualified, the policy has already been implemented and the narrative has already been captured.
\end{enumerate}
The time lag between amplification and correction is the ``laundering window''---the period during which unvetted research shapes policy through borrowed credibility.
\end{model}

\section{Crisis as Narrative Opportunity}
\label{sec:ch25-crisis}

\subsection{The shock doctrine model}

\begin{model}[Shock Doctrine]
\label{mod:shock-doctrine}
\index{shock doctrine}
Following \textcite{klein2007}, economic crises create windows of doxonomic opportunity.
During a crisis, public attention is high, uncertainty is high, and trust in existing narratives is shaken.
The first coherent narrative to fill the vacuum captures the frame.

Formally, let $H(t)$ be narrative entropy at time $t$.
A crisis causes $H(t) \to H_{\text{max}}$ (maximum uncertainty), followed by rapid entropy reduction as a new dominant frame is installed:
\[
  H(t) = H_{\text{max}} \cdot e^{-\lambda(t - t_{\text{crisis}})} + H_{\text{new}}
\]
The rate $\lambda$ depends on:
\begin{itemize}[nosep]
  \item the \emph{preparedness} of competing narratives (did someone have a pre-packaged explanation ready?),
  \item the \emph{institutional infrastructure} available to amplify the narrative,
  \item the \emph{credibility} of the narrative's proponents.
\end{itemize}
The group that supplies the new frame controls the post-crisis common sense.
\end{model}

\subsection{Narrative preparedness}

The austerity narrative's success after 2008 was not spontaneous; it reflected decades of institutional investment.
When the crisis hit, austerity proponents had:
\begin{itemize}[nosep]
  \item pre-written policy briefs and talking points,
  \item established relationships with journalists and editors,
  \item academic research (including Reinhart-Rogoff) ready for deployment,
  \item institutional credibility from decades of fiscal-conservatism advocacy,
  \item a simple, intuitive story (household metaphor) that required no economic training to understand.
\end{itemize}

Keynesian counter-narratives existed but lacked comparable infrastructure: fewer think tanks, less media presence, no simple metaphor as powerful as ``belt-tightening,'' and an argument (``spend more during a recession'') that contradicts intuitive household logic.

\section{The 2008--2013 Austerity Turn: A Detailed Case Study}
\label{sec:ch25-case}

\subsection{Phase 1: Crisis narrative (2008--2009)}

The initial frame was ``banks caused this.''
Public anger targeted financial institutions.
Narrative entropy was high.
Government intervention (bailouts, stimulus) was initially popular: the American Recovery and Reinvestment Act (February 2009, \$787 billion) passed with broad public support, and similar stimulus packages were enacted in the UK, Germany, China, and elsewhere.

\subsection{Phase 2: Reframing (2009--2010)}

Within 18 months, the dominant frame shifted from ``private-sector recklessness'' to ``public-sector profligacy.''
This reframing was not spontaneous.
It was driven by:
\begin{itemize}[nosep]
  \item \textbf{Think tanks}: the Reinhart-Rogoff paper (January 2010), the Peterson Foundation's ``fiscal responsibility'' campaign, and dozens of policy briefs from Heritage, Cato, and their European counterparts.
  \item \textbf{Financial press}: the Financial Times, Wall Street Journal, and Economist shifted framing from banking crisis to sovereign debt crisis.
  A content analysis would likely show the frequency of ``sovereign debt'' overtaking ``banking crisis'' in headlines by mid-2010.
  \item \textbf{European institutions}: the ECB, the European Commission, and the ``troika'' (ECB/EC/IMF) conditioned bailout assistance on austerity measures in Greece, Ireland, Portugal, and Spain, institutionalizing the reframing in policy.
  \item \textbf{Political actors}: the UK Conservative Party's 2010 campaign centered on ``the mess Labour left,'' reframing a global financial crisis as a domestic fiscal crisis.
  \item \textbf{Credit rating agencies}: S\&P's downgrade of US sovereign debt (August 2011) reinforced the narrative that government spending was the problem.
\end{itemize}

\subsection{Phase 3: Normalization (2010--2013)}

Austerity became common sense across much of Europe and influenced US fiscal policy (the sequester, the debt ceiling battles).
The household metaphor (``belt-tightening'') dominated media framing.
Counter-narratives (Keynesian stimulus, MMT) existed but lacked comparable institutional infrastructure.

\begin{remark}
The austerity turn illustrates several doxonomic mechanisms simultaneously:
the shock doctrine model (crisis opens the window),
narrative infrastructure effects (think tanks had pre-packaged arguments ready),
credibility laundering (an academic paper gave the 90\% threshold scientific authority),
and common-sense capture (within two years, austerity went from a policy choice to ``the only responsible option'').
\end{remark}

\subsection{Phase 4: Partial narrative erosion (2013--2020)}

The austerity consensus began eroding as:
\begin{itemize}[nosep]
  \item the Herndon correction undermined the 90\% threshold,
  \item austerity-adopting countries (Greece, UK, Spain) experienced prolonged recessions while non-austerity countries (US, with partial stimulus; Iceland, with heterodox policy) recovered faster,
  \item the IMF's own chief economist, Olivier Blanchard, published research arguing that fiscal multipliers had been underestimated---austerity was more damaging than predicted,
  \item populist movements (Syriza, Podemos, Sanders, Corbyn) explicitly challenged the austerity narrative.
\end{itemize}

But the institutional infrastructure remained: balanced-budget rules, the Maastricht criteria, credit-rating-agency discipline, and the household metaphor continued to constrain fiscal policy even as the intellectual case for austerity weakened.

\section{A Formal Model of Fiscal Belief}
\label{sec:ch25-model}

\begin{model}[Austerity Belief Dynamics]
\label{mod:austerity-dynamics}
\index{austerity belief}
Let $\belief{A}(t)$ be the strength of austerity belief in the population.
The dynamics are:
\[
  \frac{d\belief{A}}{dt} = \underbrace{\alpha \phi_A (1 - \belief{A})}_{\text{institutional production}} + \underbrace{\beta \cdot \text{Crisis}(t) \cdot (1 - \belief{A})}_{\text{crisis amplification}} - \underbrace{\gamma \cdot \Delta U(t) \cdot \belief{A}}_{\text{unemployment erosion}} - \underbrace{\delta \belief{A}}_{\text{natural decay}}
\]
where:
\begin{itemize}[nosep]
  \item $\phi_A$ is the institutional investment in austerity narratives,
  \item $\text{Crisis}(t)$ is an indicator for economic crisis periods (when the shock doctrine operates),
  \item $\Delta U(t)$ is the change in unemployment (rising unemployment associated with austerity erodes belief),
  \item $\delta$ is the natural decay rate (generational replacement, fading salience).
\end{itemize}
\end{model}

\begin{proposition}[Austerity Belief Persistence]
\label{prop:austerity-persistence}
The steady-state austerity belief (absent crisis) is:
\[
  \belief{A}^* = \frac{\alpha \phi_A}{\alpha \phi_A + \gamma \Delta U^* + \delta}
\]
If $\Delta U^* \approx 0$ (no unemployment change), the steady state depends only on the ratio of institutional production to natural decay: $\belief{A}^* = \alpha \phi_A / (\alpha \phi_A + \delta)$.
High institutional investment produces high equilibrium belief even without a crisis trigger.
\end{proposition}

\section{Distributional Analysis}
\label{sec:ch25-distributional}

\subsection{Who benefits from the austerity narrative?}

The austerity narrative's distributional consequences are highly unequal:

\begin{enumerate}
  \item \textbf{Holders of government bonds}: austerity protects the value of government debt by reducing default risk and inflation risk. Bondholders are disproportionately wealthy individuals and financial institutions.

  \item \textbf{Opponents of public spending}: austerity reduces the size of government, which aligns with the preferences of actors who view government spending as competing with private enterprise (for labor, for legitimacy, for narrative space).

  \item \textbf{Employers}: public-sector wage cuts and workforce reductions weaken the public sector as an alternative employer, increasing labor supply to the private sector and reducing bargaining power.

  \item \textbf{Tax-reduction advocates}: austerity narratives that frame spending as the problem (rather than revenue) implicitly argue for lower taxes, benefiting high-income taxpayers disproportionately.
\end{enumerate}

\begin{proposition}[Austerity as Redistribution]
\label{prop:austerity-redistribution}
\index{austerity!distributional effects}
Austerity programs redistribute income and risk from public-service users (disproportionately lower-income) to bondholders and taxpayers (disproportionately higher-income).
The fiscal morality narrative (Definition~\ref{def:fiscal-morality}) performs the doxonomic function of framing this redistribution as a moral imperative rather than a distributional choice.
\end{proposition}

\subsection{The Greek case}

\begin{example}
\label{ex:greece}
Greece's austerity program (2010--2018) provides a stark illustration:
\begin{itemize}[nosep]
  \item GDP contracted by 26\% (a depression comparable in scale to the 1930s US experience).
  \item Unemployment peaked at 27\% (youth unemployment at 60\%).
  \item Public health spending was cut by 40\%, leading to documented increases in infant mortality, HIV transmission, and mental health crises.
  \item Pension payments were cut by 40--50\%, pushing elderly poverty rates above 20\%.
  \item Meanwhile, the primary beneficiaries of the bailout funds were German and French banks that held Greek government bonds---roughly 90\% of bailout money went to debt servicing rather than the Greek economy.
\end{itemize}
The austerity narrative framed this as ``Greece living beyond its means'' and ``Greeks retiring too early''---a narrative shield (Definition~\ref{def:narrative-shield}) that deflected attention from the structural causes (eurozone design flaws, predatory lending, Northern European trade surpluses) to individual/national moral failings.
\end{example}

\section{Negative Cases and Counter-Narratives}
\label{sec:ch25-negative}

\subsection{Iceland: referendum against austerity}

Iceland rejected bank bailouts and austerity in 2008--2009 through a popular referendum.
Instead of socializing private banking losses, Iceland let the banks fail, imposed capital controls, and maintained social spending.
Its recovery was faster than most austerity-adopting European countries: GDP returned to pre-crisis levels by 2015, unemployment fell below 4\%, and the fiscal position recovered through growth rather than cuts.

The Icelandic case demonstrates that the austerity narrative requires specific institutional conditions to achieve capture.
Iceland's small size, homogeneous population, and strong civic traditions made the counter-narrative (``this was a banking crisis, not a spending crisis'') more credible than in larger, more institutionally fragmented countries.

\subsection{Latin American heterodoxy}

Several Latin American countries in the 2000s explicitly rejected IMF austerity programs: Argentina (after its 2001 default), Ecuador, Bolivia, and to some extent Brazil under Lula.
Their post-rejection economic performance was generally positive, undermining the claim that austerity is necessary for recovery.

These cases produced a regional counter-narrative infrastructure: ALBA (Bolivarian Alliance), the Banco del Sur, and academic networks around Latin American structuralism and dependency theory.
The counter-infrastructure was smaller than the pro-austerity infrastructure but was regionally effective.

\subsection{Modern Monetary Theory}

MMT has emerged as the most theoretically developed counter-narrative to fiscal morality:
\begin{itemize}[nosep]
  \item A sovereign currency issuer cannot ``run out'' of money in its own currency.
  \item The constraint on government spending is not financial (money) but real (inflation, resource availability).
  \item Government deficits are the private sector's surpluses (accounting identity).
  \item The household metaphor is therefore fundamentally misleading: a currency-issuing government is nothing like a household.
\end{itemize}

MMT's doxonomic challenge: despite logical coherence, it struggles against the household metaphor because it lacks an equally intuitive counter-metaphor.
``The government is the issuer of the currency, not a user of it'' is true but requires economic sophistication that ``families can't spend more than they earn'' does not.
The asymmetry in cognitive accessibility---not in evidential quality---explains the household metaphor's dominance.

\section{Media Framing of Fiscal Policy}
\label{sec:ch25-media}

\subsection{Content analysis evidence}

Studies of media framing of fiscal policy consistently find:
\begin{itemize}[nosep]
  \item deficit and debt stories outnumber stories about the benefits of public spending,
  \item the household metaphor is the dominant framing device,
  \item ``expert'' sources quoted on fiscal matters are disproportionately from institutions with pro-austerity positions,
  \item the distributional consequences of austerity are underreported relative to the aggregate fiscal position.
\end{itemize}

\begin{proposition}[Asymmetric Fiscal Framing]
\label{prop:asymmetric-framing}
\index{fiscal framing}
Media coverage of fiscal policy exhibits a systematic asymmetry: costs of government spending are framed as visible and immediate (``the deficit is \$X trillion''), while benefits are framed as diffuse and uncertain (``public investment may promote growth'').
Conversely, costs of spending cuts are framed as diffuse (``some programs may be reduced'') while fiscal benefits are framed as visible and immediate (``the deficit will shrink by \$X billion'').
This asymmetry favors the austerity narrative regardless of the actual evidence.
\end{proposition}

\section{Notes and References}
\label{sec:ch25-notes}

On austerity as ideology, see \textcite{streeck2016} and \textcite{brown2015}.
The shock doctrine is developed in \textcite{klein2007}.
The political economy of debt narratives connects to \textcite{graeber2011}.
For the history of neoliberal fiscal narratives, see \textcite{mirowski2009}, \textcite{harvey2005}, and \textcite{slobodian2018}.

On the Reinhart-Rogoff controversy, see Herndon, Ash, and Pollin (2014).
Blanchard and Leigh (2013) document the underestimation of fiscal multipliers.
On Iceland's alternative path, see Wade and Sigurgeirsd{\'o}ttir (2012).
On the Greek crisis, see Varoufakis (2017).

Modern Monetary Theory is developed in Wray (2015) and Kelton (2020).
On media framing of fiscal policy, see Blyth (2013).
System justification in the context of economic policy is analyzed in \textcite{jost2004}.
For the historical genealogy of fiscal morality, see \textcite{graeber2011} and \textcite{piketty2014}.

\section{Exercises}

\begin{exercise}
Construct a doxonomic genealogy of the phrase ``living beyond our means'' in political discourse.
\begin{enumerate}[nosep]
  \item When did it first appear in budget debates?
  \item What institutional channels popularized it?
  \item How does its frequency in media coverage correlate with economic conditions?
  \item Is the phrase used symmetrically (applied equally to tax cuts that increase deficits and spending that increases deficits)?
\end{enumerate}
\end{exercise}

\begin{exercise}
Using the shock doctrine model (Model~\ref{mod:shock-doctrine}), analyze the 2008 financial crisis and subsequent austerity turn.
\begin{enumerate}[nosep]
  \item Estimate the narrative entropy trajectory $H(t)$ qualitatively (sketch the curve from 2007 to 2013).
  \item Identify who supplied the dominant post-crisis frame and through what institutional channels.
  \item What counter-narratives were available and why did they fail to capture the frame?
  \item Apply the model to a second crisis (e.g., COVID-19). Did the shock doctrine operate similarly?
\end{enumerate}
\end{exercise}

\begin{exercise}
Design a survey to measure belief in the household-budget analogy for government spending.
\begin{enumerate}[nosep]
  \item Write five items that capture different aspects of fiscal morality belief.
  \item Predict which demographic groups will score highest and explain why (consider education, political affiliation, age, income).
  \item Include a treatment condition: half of respondents read a brief explanation of sovereign currency issuance before answering. Does it change responses?
\end{enumerate}
\end{exercise}

\begin{exercise}
Apply the academic credibility laundering model (Model~\ref{mod:credibility-laundering}) to a case other than Reinhart-Rogoff.
\begin{enumerate}[nosep]
  \item Identify the researcher, institution, finding, and amplification channel.
  \item Estimate the ``laundering window'' (time between amplification and correction).
  \item What policy decisions were made during the laundering window?
  \item Were those decisions reversed after the correction? If not, why not?
\end{enumerate}
\end{exercise}

\begin{exercise}
Conduct a content analysis of fiscal-policy coverage in two newspapers with different editorial stances.
\begin{enumerate}[nosep]
  \item Count the frequency of the household metaphor and its variants.
  \item Categorize expert sources by institutional affiliation.
  \item Assess whether costs and benefits of spending are framed symmetrically (Proposition~\ref{prop:asymmetric-framing}).
  \item What doxonomic explanation accounts for the patterns you find?
\end{enumerate}
\end{exercise}

\begin{exercise}
Compare the distributional consequences of austerity in Greece (Example~\ref{ex:greece}) with Iceland's alternative approach.
\begin{enumerate}[nosep]
  \item Tabulate key economic and social indicators for both countries (GDP change, unemployment, poverty, health outcomes) from 2008 to 2018.
  \item Who bore the costs in each case?
  \item Apply the legitimation framework (Chapter~\ref{ch:inequality-legitimacy}): why was the austerity narrative accepted in Greece but rejected in Iceland?
  \item What institutional differences explain the different narrative outcomes?
\end{enumerate}
\end{exercise}

\begin{exercise}[Computational]
Simulate the austerity belief dynamics model (Model~\ref{mod:austerity-dynamics}):
\begin{enumerate}[nosep]
  \item Set $\belief{A}(0) = 0.4$, $\alpha = 0.1$, $\phi_A = 3$, $\beta = 0.5$, $\gamma = 0.3$, $\delta = 0.05$.
  \item Simulate 100 periods with no crisis. Plot $\belief{A}(t)$.
  \item At $t = 30$, introduce a crisis ($\text{Crisis}(t) = 1$ for $t \in [30, 35]$). How does $\belief{A}$ respond?
  \item At $t = 50$, introduce rising unemployment ($\Delta U(t) = 0.5$ for $t \in [50, 70]$) as a consequence of the austerity policies adopted during the crisis. Does this erode austerity belief?
  \item Compare scenarios with different $\phi_A$ (narrative investment): how much investment is needed to sustain austerity belief despite rising unemployment?
\end{enumerate}
\end{exercise}
